\section{관련 연구}
자율 시스템의 안전 검증은 시스템을 블랙박스로 두고 실패를 유발하는 입력을 탐색하는
falsification 문제로 다뤄져 왔다. Corso 등은 최적화, 경로 계획, 강화학습, 중요도 샘플링 기반
블랙박스 안전 검증 기법을 정리하였고~\cite{corso}, Adaptive Stress Testing은 강화학습으로 가장 가능성
높은 실패 경로를 찾는다~\cite{ast}. 본 과제는 탐색기로 LLM을 사용하여, 수치 신호뿐 아니라 행동 기록과
로봇의 기능 호출 결과 같은 비정형 관측을 다음 시험 설계의 근거로 활용한다는 점이 다르다.

로봇 조작 분야에서는 robosuite~\cite{robosuite} 기반의 LIBERO가 언어 조건 조작 과제의 표준 벤치마크를 제공하며~\cite{libero},
GR00T~N1은 범용 휴머노이드·조작 VLA 기반 모델이다~\cite{groot}. 물리 시뮬레이션에는
MuJoCo~\cite{mujoco}를 사용하였다. 3D Scene Graph는 객체와 공간 관계를 계층 그래프로 표현하는
방식으로~\cite{sg}, 본 과제에서 테스트 공간을 구조화하는 입력 표현으로 사용한다.
본 보고서에 나오는 도구와 용어는 부록 F에 정리하였다.
